\begin{figure}[H]
\centering
\begingroup
\resizebox{\linewidth}{!}{%
\begin{minipage}{20cm}

\begin{tikzpicture}[atlas]
\node[inner sep=16pt]{%
\begin{minipage}{13.2cm}
\centering
{\color{atlasink}\rule{\linewidth}{0.8pt}}\\[2pt]
{\scshape\LARGE\color{atlasink} The McShane Integral}\\[2pt]
{\color{atlasink}\rule{\linewidth}{0.8pt}}\\[12pt]

\begin{tikzpicture}[atlas,scale=1.35]
  \def\Fexpr{1.6+0.6*sin(55*\x)}
  % --- faint context cells (ordinary terms of the McShane sum), no tags shown ---
  \foreach \xa in {0,1,3,4}{
     \pgfmathsetmacro\xb{\xa+1}
     \pgfmathsetmacro\xm{\xa+0.5}
     \pgfmathsetmacro\hh{1.6+0.6*sin(55*\xm)}
     \fill[atlasblue!9] (\xa,0) rectangle (\xb,\hh);
     \draw[atlasblue!32,line width=0.4pt] (\xa,0) rectangle (\xb,\hh);
  }
  % --- the spotlight cell [2,3] with an EXTERNAL tag t* = 1.5 ---
  \def\ts{1.5}
  \pgfmathsetmacro\hstar{1.6+0.6*sin(55*\ts)}     % height = f(t*)
  \fill[atlasorange!16] (2,0) rectangle (3,\hstar);
  \draw[atlasorange!75,line width=0.8pt] (2,0) rectangle (3,\hstar);
  % carry the height from the external tag over to the cell
  \draw[atlasorange!75,line width=0.7pt,dash pattern=on 2pt off 1.5pt] (\ts,0)--(\ts,\hstar);
  \draw[atlasorange,line width=0.8pt,-{Stealth}] (\ts,\hstar)--(2.98,\hstar);
  \node[atlas dot,fill=atlasorange] at (\ts,\hstar){};
  \draw[atlasorange,line width=1pt] (\ts,-0.07)--(\ts,0.07);
  % --- the curve over everything ---
  \glowstroke{atlasink}{0:5}{\Fexpr}
  % --- axes + partition node ticks ---
  \draw[atlas axes] (0,0)--(5.5,0);
  \draw[atlas axes] (0,0)--(0,2.95);
  \foreach \x in {0,1,2,3,4,5}{\draw[atlasink,line width=0.5pt] (\x,-0.06)--(\x,0.06);}
  % --- function-tier labels ---
  \node[atlas label,right] at (4.05,2.5){\scriptsize$y=f(x)$};
  \node[atlas label,above left=-2pt] at (\ts,\hstar){\scriptsize$f(t_\star)$};
  \node[atlas label,below] at (\ts,-0.07){\scriptsize$t_\star$};
  \node[atlas label,below] at (2,-0.07){\scriptsize$x_{i-1}$};
  \node[atlas label,below] at (3,-0.07){\scriptsize$x_i$};
  \node[atlas label] at (2.5,0.42){\scriptsize$f(t_\star)\,\Delta x_i$};

  % --- gauge tier (below the axis): the delta-fine containment ---
  \def\dd{1.7}
  \pgfmathsetmacro\wl{\ts-\dd}\pgfmathsetmacro\wr{\ts+\dd}
  \draw[atlasconegreen,line width=1.1pt] (\wl,-0.75)--(\wr,-0.75);
  \draw[atlasconegreen,line width=1.1pt] (\wl,-0.66)--(\wl,-0.84);
  \draw[atlasconegreen,line width=1.1pt] (\wr,-0.66)--(\wr,-0.84);
  \draw[atlasblue,line width=2.1pt] (2,-0.75)--(3,-0.75);       % the cell I_i
  \node[atlas dot,fill=atlasorange,minimum size=2.8pt] at (\ts,-0.75){}; % tag OUTSIDE the cell
  \draw[atlasgray,line width=0.4pt,dash pattern=on 1.5pt off 1.5pt] (2,0)--(2,-0.75);
  \draw[atlasgray,line width=0.4pt,dash pattern=on 1.5pt off 1.5pt] (3,0)--(3,-0.75);
  \draw[atlasgray,line width=0.4pt,dash pattern=on 1.5pt off 1.5pt] (\ts,0)--(\ts,-0.75);
  \node[atlas label,left] at (\wl,-0.75){\scriptsize$t_\star-\delta(t_\star)$};
  \node[atlas label,right] at (\wr,-0.75){\scriptsize$t_\star+\delta(t_\star)$};
  \node[atlas label,below] at (\ts,-0.88){\scriptsize$t_\star\notin[x_{i-1},x_i]$};
  \node[atlas label,fill=atlasconegreen!18,inner sep=2pt,rounded corners=1pt] at (3.6,-0.75)
     {\scriptsize $I_i\subseteq\big(t_\star-\delta(t_\star),\,t_\star+\delta(t_\star)\big)$};
\end{tikzpicture}

\vspace{8pt}
{\scshape\footnotesize\color{atlasink} Figure. The McShane Integral --- the untethered tag: a $\delta$-fine partition whose tags roam free.}\\[6pt]

% --- formal strip + positioning ---
\footnotesize\color{atlasink}
\begin{minipage}{12.4cm}
\setlength{\parskip}{3pt}
\textbf{Rule (the only change from Henstock--Kurzweil).} A tagged partition is
\emph{McShane $\delta$-fine} iff
$\ \forall i\ \big(I_i\subseteq(t_i-\delta(t_i),\,t_i+\delta(t_i))\big)$,
with tags $t_i\in[a,b]$ \emph{arbitrary} --- the clause $t_i\in I_i$ is dropped.
The sum is still $S=\sum_i f(t_i)\,\Delta x_i$, and
\[
  \textstyle f\ \text{McShane integrable}\ \Longleftrightarrow\ f\in L^1[a,b],
  \qquad\text{with equal values.}
\]
\textbf{Where it sits:}\quad
$\mathcal R\ \subsetneq\ \mathcal L=\text{McShane}\ \subsetneq\ \mathcal{HK}$
\ (Riemann $\subsetneq$ Lebesgue $=$ McShane $\subsetneq$ Henstock--Kurzweil).
\end{minipage}
\end{minipage}%
};
\end{tikzpicture}
\end{minipage}%
}
\endgroup
\caption{The McShane integral and the untethered tag}
\label{fig:mcshane}
\end{figure}
